\section{Strict realization of the coherent tail presentation and repair in degrees zero and one}
\label{sec:136-coherent-pd-realization}

\paragraph{Archive scope and inputs.}
The coefficient and vertical requirements are cote136, physical
p.~117, footer116, scan117; the horizontal requirement is physical
p.~119, footer118, scan119; the common-source question is physical
p.~121, footer120, scan121. The two inputs used here are the fully
checked ARCH-136-RELATIVE-COHERENT-MOORE and the supplied
AQ-136-RELATIVE-H1. We construct an actual strict PD extension
realizing the former and verify the hypotheses of the latter.
No integral Poincare lemma for a nonflat parameter module is used.
The resulting statement controls form degrees zero and one; it
does not control the remaining integral cohomology in form degree
two or prove the common-source assertion.

Put \(S=\Gamma_{\mathbb Z}(T)\). Let \(A\) be a simplicial,
nonnegatively bigraded strict PD differential \(S\)-algebra,
with \(A_0=S\) and \(A_n^{0,*}=\mathbb Z[0]\), equipped with
strict maps \(f_j:A\to C_j\) to a fixed finite family of
admissible theories. Suppose, for one bound \(W\) independent
of the simplex, that
\[
 H^0((A_{\mathbb Q})_n^{w,*})=S_{\mathbb Q}^w,
 \qquad H^1((A_{\mathbb Q})_n^{w,*})=0
                       \quad(w\geq W).                 \tag{136.CP1}
\]
The coefficient identification is the canonical one.

Choose an odd integer \(q\geq1\), an integer \(p\geq2\),
a weight \(m>0\), a simplicial pair \(Y\subseteq X\)
containing every vertex of \(X\), and
\[
 a\in A^{m,q}(X,Y),\quad b=da,\quad
 z=\gamma_p(b),\quad e=a\gamma_{p-1}(b),\quad
 D=p(q+1),\quad v=pm.                                  \tag{136.CP2}
\]
Thus \(D\) is even, \(dz=0\), and \(de=pz\).
Assume that every input row \((C_j)_*^{m,q}\) is horizontally
acyclic, including in degree zero. Choose \(M\) sufficiently
large for ARCH-136-RELATIVE-COHERENT-MOORE to apply to all the
targets, and choose its weight-graded free resolution
\[
 \cdots\longrightarrow P_2\xrightarrow{\partial_2}P_1
 \xrightarrow{\partial_1}P_0\xrightarrow{\epsilon}I_M
 \longrightarrow0,\qquad
 I_M=\bigoplus_{w\geq M}\mathbb ZT^{[w]}.                \tag{136.CP3}
\]
All generators of every \(P_i\) have weight at least \(M\).
Use one resolution for all targets. The cited theorem gives
relative target maps \(U_i^{(j)},V_i^{(j)}\) obeying its full
relations, including the relation in form degree zero.

\paragraph{ARCH-136-COHERENT-PD-REALIZATION.}
There is an actual simplicial strict PD differential algebra
\(E\), with maps
\[
                  A\longrightarrow E\longrightarrow C_j,
                                                               \tag{136.CP4}
\]
extending every \(f_j\), and with \(E_0=S\) and
\(E_n^{0,*}=\mathbb Z[0]\). Every coefficient-tail multiple
\(\epsilon(u)z\), for \(u\in P_0\), has its specified
primitive in \(E(X,Y)\). At each simplex, the map
\(A_{\mathbb Q}\to E_{\mathbb Q}\) is a cochain homotopy
equivalence preserving weights. More precisely, after an
explicit triangular change of generators it is the inclusion
into a free tensor product of rational unit disks.

In every simplicial degree, the new module of generators in
form degree zero is weightwise free as an abelian group. This
does not assert that its divided-power algebra is torsion-free
or that \(A\to E\) is an integral cohomology equivalence.

\begin{proof}
Let \(L=\mathbb Z[X]/\mathbb Z[Y]\), a simplicial abelian
group with \(L_0=0\), whose terms are free abelian on the
simplexes outside \(Y\). Tensor each parameter module in
(136.CP3) with \(L\), and shift its weights by \(v\).
Introduce \(S\)-linear generators
\[
 \begin{array}{lll}
 U_i:P_i[v]\otimes L\longrightarrow E^{*,D-1-i},
       &&0\leq i\leq D-1,\\
 V_i:P_i[v]\otimes L\longrightarrow E^{*,D-2-i},
       &&0\leq i\leq D-2.
 \end{array}                                             \tag{136.CP5}
\]
The free graded extension uses divided powers on even generators
and strictly exterior odd generators. Its differential on these
generators is prescribed by
\[
 \begin{aligned}
 dU_0&=f_z,&
 dU_i&=-U_{i-1}\partial_i\quad(i\geq1),\\
 dV_0&=pU_0-f_e,&
 dV_i&=pU_i+V_{i-1}\partial_i\quad(i\geq1),
 \end{aligned}                                             \tag{136.CP6}
\]
where \(f_z(u\otimes\sigma)=\epsilon(u)z_\sigma\), and
similarly for \(f_e\). The relative condition makes these
well-defined on \(L\). All formulas are simplicial and
weight-preserving. Impose, as an identifying PD ideal, the
following form-zero relations:
\[
 U_{D-1}\partial_D=0,
 \qquad pU_{D-1}+V_{D-2}\partial_{D-1}=0.                \tag{136.CP7}
\]
These relations include the entire bottom relation, rather
than just the factorization through the last syzygy.

We check existence of the differential before invoking any
cohomology claim. The identities \(df_e=pf_z\),
\(f_e\partial_1=f_z\partial_1=0\), and
\(\partial_{i-1}\partial_i=0\) show directly that the square
of (136.CP6) vanishes on every generator. The differential of
the first relation in (136.CP7) is zero by \(\partial^2=0\).
The differential of the second is
\[
 -pU_{D-2}\partial_{D-1}
 +(pU_{D-2}+V_{D-3}\partial_{D-2})\partial_{D-1}=0.
\]
There is no missing small-index case, since \(D\geq4\).
The special divided-power rule makes the full identifying PD
ideal differential-stable. The square of the differential on
even divided powers vanishes because the differential of an
even element is odd and strictly squares to zero. This
constructs \(E\) as a strict PD differential algebra.

The coherent target maps supplied by
ARCH-136-RELATIVE-COHERENT-MOORE obey exactly (136.CP6)--(136.CP7),
so the universal property gives every map in (136.CP4).
No target flatness is required. Because \(L_0=0\), no new
generator occurs in simplicial degree zero. Every new generator
has positive weight. Hence the assertions about degree zero
and weight zero follow. In particular coefficients stay
injective in every simplex: their degeneracy maps from
\(E_0=S\) have face retractions.

For clarity we compute the bottom parameter module integrally.
Fix a simplex, suppress the harmless tensor factor \(L_n\)
and the weight shift, and put
\[
 P=P_{D-2},\qquad
 Q=\operatorname{coker}(P_D\to P_{D-1})
       \xrightarrow{\ \partial_{D-1}\ }P.
                                                               \tag{136.CP8}
\]
The inclusion follows from exactness of (136.CP3). The form-zero
parameter module is
\[
 M_0=(P\oplus Q)/\langle(q,pq):q\in Q\rangle
          \ \cong\ P+p^{-1}Q\ \subseteq P[1/p],          \tag{136.CP9}
\]
where \((x,y)\mapsto x-y/p\). To prove injectivity, note
that \(P/Q\) embeds into the free \(S\)-module
\(P_{D-3}\); it is therefore torsion-free over \(\mathbb Z\).
If \(px=y\in Q\), this torsion-freeness gives \(x\in Q\),
and \((x,y)=(x,px)\) is exactly one of the indicated relations.
Surjectivity onto the displayed lattice is immediate.
Weightwise \(P/Q\) is a subgroup of a free abelian group and
is therefore free abelian, with no finite-rank assumption.
The short exact sequence \(0\to Q\to P\to P/Q\to0\)
splits as abelian groups in each weight. Consequently
\(P+p^{-1}Q\) is weightwise free abelian as well. Tensoring
with the free group \(L_n\) gives the same conclusion for
all the original parameters. This is a statement about the
linear parameter module only.

Now extend scalars to \(\mathbb Q\). The second relation
in (136.CP7) eliminates \(U_{D-1}\); its first relation is
then automatic. For \(0\leq i\leq D-2\), replace \(U_i\)
by the following coordinates on its free parameter module:
\[
 Z_0=pU_0-f_e,\qquad
 Z_i=pU_i+V_{i-1}\partial_i\quad(i\geq1).               \tag{136.CP10}
\]
This is an invertible triangular change over the old algebra
\(A_{\mathbb Q}\): the inverse divides by \(p\) and uses
only the displayed old elements and the \(V\)'s. The
relations and (136.CP6) now read
\[
                      dV_i=Z_i,\qquad dZ_i=0.
                                                               \tag{136.CP11}
\]
There are no remaining relations among these free parameters.
In characteristic zero the strict PD algebra is the usual free
graded-commutative algebra with its determined divided powers.
Thus (136.CP11) expresses \(E_{\mathbb Q}\), at each simplex,
as \(A_{\mathbb Q}\) tensored over \(S_{\mathbb Q}\) with
the free unit disks on bases of the free modules
\(P_i[v]\otimes L_n\), for \(0\leq i\leq D-2\).

Here is a contraction which also handles arbitrary parameter
ranks. Give each new coordinate \(V_i,Z_i\) arity one and
old elements arity zero. On the rational free algebra let
\(h\) be the odd derivation with
\(h(Z_i)=V_i\), \(h(V_i)=0\), and \(h(A_{\mathbb Q})=0\).
The graded derivation calculation gives
\(dh+hd=r\) on arity \(r\). Dividing by \(r>0\) contracts
the positive-arity part; arity zero is \(A_{\mathbb Q}\).
Every polynomial has finite arity and finite support, so this
defines a genuine contraction for arbitrary ranks. It
preserves the original weight, and the augmentation setting
all the new coordinates to zero is a differential algebra
retraction. This proves the claimed rational cochain homotopy
equivalence and every assertion of the statement.
\end{proof}

\paragraph{ARCH-136-COHERENT-LOW-REPAIR.}
For the same data, there is a further sequence of strict maps
under \(E\) and over all targets, with colimit \(E^+\), such
that, for one \(W^*\) independent of the simplex,
\[
 E^+_0=S,\quad (E^+)^{0,*}=\mathbb Z[0],\quad
 H^0((E^+)_n^{w,*})=S^w,\quad
 H^1((E^+)_n^{w,*})=0\quad(w\geq W^*).                  \tag{136.CP12}
\]
All the primitives installed in (136.CP6) persist. The bound
\(W^*\) can be taken to be the maximum of the bound in
(136.CP1) and the vertical degree-zero/degree-one bounds of
the finitely many targets. No integral flatness assertion is
part of this conclusion.

\begin{proof}
The rational cochain homotopy equivalence just proved shows
that \(E\) satisfies (136.CP1) with the same bound as \(A\).
The other hypotheses of AQ-136-RELATIVE-H1 are \(E_0=S\),
weight-zero equality, and the already constructed strict maps
to the admissible targets. All have been verified. Apply that
theorem with the stated maximum bound. It supplies the actual
strict PD maps and filtered colimit (136.CP12). In particular
this application is justified by the explicit rational
coordinate calculation; it does not import an acyclicity
assertion for a nonflat even disk. Identities such as
\(dU_0(u)=\epsilon(u)z\) are preserved by every differential
algebra map and hence remain true in the colimit.
\end{proof}

\paragraph{Horizontal bounds and the exact remaining issue.}
Neither the free PD quotient defining \(E\) nor
AQ-136-RELATIVE-H1 individually preserves horizontal acyclicity.
If a horizontally acyclic source is needed, the accepted
AQ-136-PARTIAL-SOURCE alternation can be applied after this
step: fix once and for all row cutoffs above the target row
bounds, attach all required free horizontal unit disks, apply
the same degree-zero/degree-one repair, and repeat. Over
\(\mathbb Q\) a free horizontal unit disk is vertically
contractible, so (136.CP1) survives each such step with the
same bound. Target row acyclicity supplies the normalized-disk
lifts. Every horizontal cycle appears at a finite stage and
is killed at a subsequent horizontal stage; every unwanted
vertical class in degree zero or one is treated at a
subsequent low-degree stage. The cutoff for each fixed row
therefore remains fixed throughout. The filtered colimit has
all horizontal tails and (136.CP12), and retains every installed
coherent primitive and every target map. This paragraph is
an application of the already supplied alternation, not a new
claim of integral vertical exactness.

The nonflat even-disk examples explain why the intervening
low-degree repair cannot simply be omitted. They do not block
the construction above. Conversely, this construction gives
no bound forcing the newly introduced integral cohomology in
degrees at least two to vanish. It reaches form degree zero
for every \(p\), and it supplies no lower form cutoff tending
to infinity with \(p\). After iteration, new cycles need not
have the native derivative form \(\gamma_p(da)\) to which
ARCH-136-RELATIVE-COHERENT-MOORE applies. These are the remaining
issues for using this realization in an integral admissible
common-source construction.
